\chapter{Introducción a la Ingeniería del Software}

\section{El producto software, propiedades y ciclo de vida}

El software no se parece a los productos que fabrica el resto de las ingenierías.
No se gasta con el uso, no está sujeto a tolerancias de material y su coste de
replicación es prácticamente nulo: una copia vale lo mismo que un millón. Lo que
cuesta es concebirlo y mantenerlo, y ahí es donde se concentra el esfuerzo de la
disciplina.

De esa naturaleza salen dos consecuencias que condicionan todo lo demás. La
primera es que no hay límites físicos que frenen la complejidad. Un puente no
puede crecer sin que la estructura lo acuse; un sistema software sí, y por eso
puede llegar a un tamaño que ninguna persona abarca. La segunda es que el
software no se deteriora, pero sí se degrada: cada cambio que se le hace lo aleja
del diseño con el que nació, y la facilidad para modificarlo acaba siendo la causa
de que se vuelva difícil de modificar \cite{sommerville2016}.

Conviene separar dos términos que se usan como sinónimos y no lo son. El
\textbf{producto software} es lo que se entrega: los programas, los datos de
configuración y la documentación que permite instalarlo, usarlo y mantenerlo. El
\textbf{sistema software} incluye además el entorno en el que ese producto opera.
Entregar solo los ejecutables no es entregar el producto.

\subsection*{Propiedades de un producto software}

La calidad de un producto software no se mide únicamente por que haga lo que se le
pide. Las propiedades que suelen exigirse son:

\begin{itemize}
    \item \textbf{Corrección y fiabilidad.} Que haga lo especificado y que siga
          haciéndolo en condiciones adversas, incluidas las que no se previeron.
    \item \textbf{Mantenibilidad.} Que pueda evolucionar con un coste razonable.
          Es la propiedad que más peso tiene a largo plazo, porque el grueso del
          gasto de un sistema se produce después de la primera entrega.
    \item \textbf{Eficiencia.} Que use de forma razonable el tiempo de cómputo, la
          memoria y el resto de recursos.
    \item \textbf{Usabilidad.} Que las personas a las que va dirigido puedan
          usarlo sin formación desproporcionada.
    \item \textbf{Seguridad y protección.} Que no cause daños y que resista un uso
          malintencionado.
\end{itemize}

Estas propiedades compiten entre sí. Un diseño muy eficiente suele ser menos
legible y, por tanto, menos mantenible; una interfaz muy flexible complica la
validación. Parte del trabajo de ingeniería consiste precisamente en decidir qué
se sacrifica y dejar constancia de por qué.

\subsection*{El ciclo de vida}

El ciclo de vida es el conjunto de etapas por las que pasa un producto desde que
se identifica la necesidad hasta que se retira. Con distintos nombres, casi todos
los modelos reconocen las mismas actividades: especificación, diseño e
implementación, validación y evolución.

Lo que distingue a un modelo de otro no son las actividades, sino cómo se
ordenan y con qué frecuencia se repiten. Merece la pena insistir en la última
etapa: \textbf{la evolución no es un apéndice del ciclo de vida, es la mayor parte
de él}. Un sistema en producción recibe cambios durante años, y el diseño con el
que se construyó determina lo caros que resultan.

\section{Concepto de Ingeniería del Software}

El término nació en la conferencia de la OTAN de 1968, convocada porque los
proyectos software se entregaban tarde, por encima del presupuesto y sin la
calidad prometida. La expresión «crisis del software» describía esa situación, y
la propuesta fue tratar el desarrollo como una ingeniería: con métodos, medidas y
procesos, en lugar de como un oficio artesanal.

La Ingeniería del Software es, en esa línea, la disciplina que se ocupa de todos
los aspectos de la producción de software, desde las primeras etapas de la
especificación hasta el mantenimiento posterior a la entrega
\cite{sommerville2016}. Dos matices de esa definición importan:

\begin{itemize}
    \item \textbf{Todos los aspectos}, no solo la programación. La gestión del
          proyecto, la elección de herramientas y la organización del equipo
          forman parte de la disciplina.
    \item \textbf{Hasta el mantenimiento}, no hasta la entrega. Un método que
          funciona bien hasta el día del despliegue y deja el sistema
          intratable después no ha resuelto el problema.
\end{itemize}

Conviene distinguirla de dos disciplinas vecinas. La \textbf{informática} estudia
los fundamentos teóricos del cómputo; la ingeniería del software se ocupa de
construir sistemas útiles con esos fundamentos, aceptando restricciones de plazo,
coste y personas. La \textbf{ingeniería de sistemas} es más amplia: abarca el
hardware, los procesos y la organización, y el software es solo una de sus
partes.

\subsection*{Responsabilidad profesional}

La disciplina lleva asociada una dimensión ética que no es decorativa. Quien
construye software decide, en la práctica, sobre la privacidad de terceros, sobre
la accesibilidad de un servicio y, en sistemas críticos, sobre la seguridad de
las personas. Los códigos deontológicos de ACM e IEEE recogen los compromisos
habituales: confidencialidad, competencia declarada con honestidad, respeto a la
propiedad intelectual y no aceptar encargos que se sepan perjudiciales
\cite{sommerville2016}.

\section{El proceso de desarrollo de software}

Un proceso de desarrollo es un conjunto ordenado de actividades que produce un
producto software. Los modelos que se estudian a continuación no son recetas
excluyentes: en un proyecto real se combinan, y la elección depende de cuánto se
conoce el problema al empezar y de cuánto cuesta equivocarse.

\subsection*{Modelo en cascada}

Ordena las actividades en fases sucesivas —especificación, diseño,
implementación, verificación y mantenimiento— y exige cerrar cada una antes de
abrir la siguiente. Su ventaja es la trazabilidad: cada fase deja documentos
revisables, lo que encaja bien en desarrollos con requisitos contractuales o
sujetos a certificación.

Su inconveniente es el que lo hizo célebre: los requisitos rara vez se conocen por
completo al principio, y el modelo hace caro cambiarlos. Un error de
especificación no se detecta hasta la validación, cuando corregirlo cuesta
órdenes de magnitud más que en su fase de origen \cite{pressman2021}.

\subsection*{Desarrollo incremental}

Divide el sistema en incrementos que se especifican, construyen y entregan por
separado, empezando por los de mayor valor. Cada entrega es funcional y sirve para
obtener realimentación antes de comprometerse con el resto.

Reduce el coste de acomodar cambios y adelanta el momento en que el cliente ve
algo que funciona. A cambio, el proceso es menos visible desde fuera —no hay
documentos de fase que marquen el avance— y la estructura del sistema tiende a
degradarse si no se reserva esfuerzo para reorganizarla.

\subsection*{Modelos evolutivos y por prototipos}

Cuando el problema no está claro, resulta razonable construir una versión reducida
que sirva para entenderlo. El prototipo puede ser \emph{desechable}, si su único
fin es aclarar requisitos y luego se descarta, o \emph{evolutivo}, si se refina
hasta convertirse en el producto. El riesgo del segundo es conocido: un prototipo
se construye sin las garantías que se exigen a un producto, y promoverlo tal cual
deja una base frágil.

El \textbf{modelo en espiral} de Boehm organiza el desarrollo en ciclos y sitúa el
análisis de riesgos como actividad explícita al principio de cada uno. Su
aportación duradera no es la forma de la espiral, sino la idea de que el proceso
debe adaptarse al riesgo asumido.

\subsection*{Métodos ágiles}

Los métodos ágiles parten de una constatación: en muchos proyectos el coste de
predecir es mayor que el de adaptarse. Priorizan las entregas frecuentes, la
colaboración con el cliente y la respuesta al cambio por encima de la planificación
detallada y la documentación exhaustiva. Scrum organiza el trabajo en iteraciones
cortas con roles y reuniones definidos; la programación extrema añade prácticas
técnicas como la integración continua, la programación por parejas y el desarrollo
guiado por pruebas \cite{beck2002,lasa2017}.

No son la respuesta a todo. Su eficacia se apoya en equipos pequeños, en un cliente
disponible y en un dominio donde equivocarse sea barato. En sistemas con
certificación obligatoria, con muchos equipos coordinados o con un coste de fallo
muy alto, conviene un proceso más ceremonioso, aunque adopte prácticas ágiles en
lo técnico.
